\subsection{Power Series}

% **Role:** convergence of functions and analytic approximation.
%
% Power series, Taylor series, analytic functions, Weierstrass approximation, Stone--Weierstrass.


Power series are controlled infinite polynomials, and approximation is uniform closure. This subsection covers radii, Taylor structure, Abel behavior, Weierstrass, Stone--Weierstrass, and elementary analytic functions.

\begin{definition} \label{def:power-approx-fourier-core-source-todo}\label{def:power-approx-core-objects} \hlt{Core Objects: Power Series, Approximation, and Gamma}
\begin{enumerate}[label=(\roman*)]
\setlength{\itemsep}{0pt}
\item A real \hlt{power series about $a$} is $\sum\limits_{n=0}^{\infty}c_n(x-a)^n$; a local representation of $f$ by one makes $f$ \hlt{analytic} at $a$.
\item A \hlt{polynomial approximant} to $f$ on a compact set $K$ is a polynomial $P$ chosen so that $\norm{f-P}_{\infty,K}$ is small.
\item For $x>0$, the \hlt{gamma function} is $\Gamma(x)\defeq\int_0^\infty t^{x-1}e^{-t}\,dt$ when the improper integral converges.
\end{enumerate}
\hlt{Key Insight}: Power series are local analytic objects; approximation theorems turn algebraic families into dense function families; $\Gamma$ extends factorials by integration.
\end{definition}

\begin{definition} \label{def:power-series-radius-basic} \hlt{Power Series and Radius of Convergence}\\
Let $(c_n)_{n\geq0}\subseteq\C$ and $\alpha\defeq\limsup\limits_{n\to\infty}\sqrt[n]{\abs{c_n}}\in[0,+\infty]$.
\begin{enumerate}[label=(\roman*)]
\setlength{\itemsep}{0pt}
\item The series $\sum\limits_{n=0}^{\infty}c_nz^n$ is a \hlt{power series}.
\item The \hlt{radius of convergence} is $R\defeq\alpha^{-1}\in[0,+\infty]$, with $0^{-1}=+\infty$ and $(+\infty)^{-1}=0$.
\end{enumerate}
\end{definition}

\begin{theorem} \label{thm:power-series-radius-convergence} \hlt{Convergence Inside the Circle}\\
For $\sum c_nz^n$, radius $R$: converge if $\abs{z}<R$, diverge if $\abs{z}>R$; boundary $\abs{z}=R$ is undecided.
\begin{proof}
For $a_n=c_nz^n$, $\limsup\limits_{n\to\infty}\sqrt[n]{\abs{a_n}}=\abs{z}\limsup\limits_{n\to\infty}\sqrt[n]{\abs{c_n}}=\frac{\abs{z}}{R}$, with extended meaning if $R=0$ or $R=+\infty$.\\
If $\abs{z}<R$, this value is $<1$; if $\abs{z}>R$, it is $>1$. Apply Theorem \ref{thm:root-test}.\\
\hlt{Structural Commentary}: Radius controls the interior and exterior, not the boundary.\\
\hlt{Key Insight}: The root test is the radius-of-convergence theorem in disguise.
\end{proof}
\end{theorem}

\begin{theorem} \label{thm:power-series-boundary-dirichlet} \hlt{A Boundary Test for Power Series}\\
If $\sum c_nz^n$ has radius $1$, $c_n\downarrow0$, and $c_n\geq0$, then it converges for $\abs{z}=1$, except possibly at $z=1$.
\begin{proof}
For $\abs{z}=1$ and $z\neq1$, put $a_n=z^n$ and $b_n=c_n$.\\
The partial sums of $\sum z^n$ are bounded because $\abs{\sum\limits_{m=0}^{n}z^m}=\abs{\frac{1-z^{n+1}}{1-z}}\leq\frac{2}{\abs{1-z}}$. Theorem \ref{thm:dirichlet-test-series} applies.\\
\hlt{Structural Commentary}: Boundary convergence away from $z=1$ comes from bounded geometric partial sums.\\
\hlt{Key Insight}: Monotone coefficients can give conditional boundary convergence.
\end{proof}
\end{theorem}

\begin{theorem} \label{thm:abel-endpoint-theorem} \hlt{Abel's Endpoint Theorem}\\
Let $s_n\defeq\sum_{k=0}^{n}c_k\to s$. For $\abs{x}<1$, set $f(x)\defeq\sum\limits_{n\geq0}c_nx^n$. Then $\lim\limits_{x\to1^-}f(x)=s$.
\begin{proof} Set $s_{-1}\defeq0$. Since $c_n=s_n-s_{n-1}$, for $\abs{x}<1$, $f(x)=\sum\limits_{n\geq0}(s_n-s_{n-1})x^n=(1-x)\sum\limits_{n\geq0}s_nx^n$.\\
Thus $f(x)-s=(1-x)\sum\limits_{n\geq0}(s_n-s)x^n$. Given $\varepsilon>0$, choose $N$ with $\abs{s_n-s}<\varepsilon/2$ for $n>N$.\\
For $0<x<1$, $\abs{f(x)-s}\leq(1-x)\sum\limits_{n=0}^{N}\abs{s_n-s}x^n+\frac{\varepsilon}{2}(1-x)\sum\limits_{n>N}x^n$.\\
The first term tends to $0$ as $x\to1^-$, and the second is $\leq\varepsilon/2$; hence $f(x)\to s$.\\
\hlt{Structural Commentary}: Endpoint convergence is one-sided Abel continuity.\\
\hlt{Key Insight}: $(1-x)\sum s_nx^n$ is a geometric average of partial sums.
\end{proof}
\end{theorem}

\begin{theorem} \label{thm:abel-product-theorem} \hlt{Abel Product Theorem for Cauchy Products}\\
Let $\sum a_n=A$ and $\sum b_n=B$. Define $c_n\defeq\sum\limits_{k=0}^{n}a_kb_{n-k}$, and suppose $\sum c_n=C$. Then $C=AB$.
\begin{proof}
For $\abs{x}<1$, set $f(x)\defeq\sum\limits_{n\geq0}a_nx^n$, $g(x)\defeq\sum\limits_{n\geq0}b_nx^n$, and $h(x)\defeq\sum\limits_{n\geq0}c_nx^n$. Since $(a_n),(b_n),(c_n)$ are bounded, these series converge absolutely on $\abs{x}<1$. Hence $f(x)g(x)=\sum\limits_{n\geq0}\bigl(\sum\limits_{k=0}^{n}a_kb_{n-k}\bigr)x^n=\sum\limits_{n\geq0}c_nx^n=h(x)$.\\
By Abel's endpoint thm, $f(x)\to A$, $g(x)\to B$, $h(x)\to C$ as $x\to1^-$;$C=\lim\limits_{x\to1^-}h(x)=\lim\limits_{x\to1^-}f(x)g(x)=AB$.\\
\hlt{Structural Commentary}: Regularize by $x$, multiply inside, pass to $x=1$.\\
\hlt{Key Insight}: Abel moves the safe product identity to the endpoint.
\end{proof}
\end{theorem}

\begin{perspective} \label{per:centered-power-series-translation} \hlt{Centered Series Are Translated Series}\\
$\sum c_n(x-a)^n$ is $\sum c_ny^n$ with $y=x-a$ and radius $R$: absolute convergence for $\abs{x-a}<R$, divergence for $\abs{x-a}>R$, and separate endpoint tests for $\abs{x-a}=R$.
\end{perspective}

\begin{figure}[H]
  \centering
  \begin{tikzpicture}[
  >=Stealth,
  every node/.style={font=\scriptsize},
  axis/.style={->, gray!70, line width=0.5pt},
  conv/.style={blue!70!black, line width=1.35pt},
  inner/.style={orange!85!black, line width=1pt, line cap=round},
  outside/.style={red!65!black, line width=0.85pt},
  boundary/.style={circle, draw=blue!70!black, fill=white, inner sep=1.6pt},
  point/.style={circle, fill=blue!70!black, inner sep=1.5pt},
  guide/.style={gray!45, dashed, line width=0.4pt}
]
  \draw[axis] (-4.3,0) -- (4.4,0) node[right] {$x$};

  \draw[outside] (-4.05,0) -- (-2.25,0);
  \draw[outside] (2.25,0) -- (4.05,0);
  \draw[conv] (-2.2,0) -- (2.2,0);
  \draw[inner] (-1.35,0.48) -- (1.35,0.48);
  \draw[inner] (-1.35,0.38) -- (-1.35,0.58);
  \draw[inner] (1.35,0.38) -- (1.35,0.58);

  \node[boundary] at (-2.2,0) {};
  \node[boundary] at (2.2,0) {};
  \node[point] at (0,0) {};

  \draw[guide] (-2.2,0.08) -- (-2.2,0.58);
  \draw[guide] (2.2,0.08) -- (2.2,0.58);

  \node[below] at (-2.2,-0.05) {$a-R$};
  \node[below] at (0,-0.05) {$a$};
  \node[below] at (2.2,-0.05) {$a+R$};
  \node[above] at (0,0.62) {$[a-\rho,a+\rho]$};
  \node[blue!65!black, below] at (0,-0.42) {$\abs{x-a}<R$};
\end{tikzpicture}

\caption{Inside $\abs{x-a}<R$: absolute and locally uniform; outside divergence; endpoints undecided.}
  \label{fig:power-series-radius-interval}
\end{figure}

\begin{definition} \label{def:real-analytic-at-point} \hlt{Real Analytic at a Point}\\
Let $E\subseteq\R$, $f:E\to\R$, $a\in\operatorname{int}E$. The function $f$ is \hlt{real analytic at $a$} if $\exists r>0$ and $(c_n)_{n\geq0}\subseteq\R$ such that:
\begin{enumerate}[label=(\roman*)]
\setlength{\itemsep}{0pt}
\item $(a-r,a+r)\subseteq E$.
\item The power series $\sum c_n(x-a)^n$ has radius of convergence $\geq r$.
\item $f(x)=\sum\limits_{n=0}^{\infty}c_n(x-a)^n$ for all $x\in(a-r,a+r)$.
\end{enumerate}
If $E$ is open and $f$ is real analytic at every $a\in E$, then $f$ is \hlt{real analytic on $E$}.
Analytic functions form the class $C^\omega$.
\end{definition}

\begin{theorem} \label{thm:power-series-uniform-differentiation} \hlt{Power Series May Be Differentiated Term by Term}\\
Let $f(x)=\sum_{n\ge0}c_nx^n$ have radius $R>0$, and set $I_R\defeq(-R,R)$, $K_\rho\defeq[-\rho,\rho]$ for $0<\rho<R$.
\begin{enumerate}[label=(\roman*)]
\setlength{\itemsep}{0pt}
\item $\sum c_nx^n$ converges uniformly on each $K_\rho$, and $f\in C^\infty(I_R)$.
\item For $k\geq1$ and $x\in I_R$, $f^{(k)}(x)=\sum\limits_{n\ge k}\frac{n!}{(n-k)!}c_nx^{n-k}$; in particular, $f^{(k)}(0)=k!c_k$.
\end{enumerate}
\begin{proof}
\begin{enumerate}[label=(\roman*)]
\setlength{\itemsep}{0pt}
\item If $0<\rho<R$, then $\sum\abs{c_n}\rho^n<\infty$; hence $\abs{c_nx^n}\leq\abs{c_n}\rho^n$ on $K_\rho$, so $\sum c_nx^n$ converges uniformly there.
\item For fixed $k$, $\frac{n!}{(n-k)!}\sim n^k$, so $\limsup\sqrt[n]{\abs{\frac{n!}{(n-k)!}c_n}}=\limsup\sqrt[n]{\abs{c_n}}$.\\
Hence every derivative series has radius $R$ and converges uniformly on each $K_\rho$.\\
Iterating termwise differentiation gives the formula and $f^{(k)}(0)=k!c_k$.
\end{enumerate}
\hlt{Structural Commentary}: Work on $K_\rho\Subset I_R$ to justify termwise operations.\\
\hlt{Key Insight}: Polynomial coefficient growth does not change the radius.
\end{proof}
\end{theorem}

\begin{corollary} \label{cor:centered-power-series-termwise-calculus} \hlt{Termwise Calculus for Centered Power Series}\\
Let $f(x)=\sum\limits_{n=0}^{\infty}c_n(x-a)^n$ have radius $R>0$.
\begin{enumerate}[label=(\roman*)]
\setlength{\itemsep}{0pt}
\item The series converges locally uniformly on $(a-R,a+R)$ and defines a continuous function.
\item For $\abs{x-a}<R$, $f'(x)=\sum\limits_{n=1}^{\infty}n c_n(x-a)^{n-1}$, locally uniformly.
\item If $[y,z]\subset(a-R,a+R)$, then $\int_y^z f(x)\,dx=\sum\limits_{n=0}^{\infty}c_n\frac{(z-a)^{n+1}-(y-a)^{n+1}}{n+1}$.
\end{enumerate}
\begin{proof}
Set $g(u)\defeq\sum c_nu^n$ and $T(x)\defeq x-a$, so $f=g\circ T$ on $(a-R,a+R)$.
\begin{enumerate}[label=(\roman*)]
\setlength{\itemsep}{0pt}
\item Local uniform convergence of $\sum c_nu^n$ on $(-R,R)$ implies local uniform convergence of $\sum c_n(Tx)^n$ on $(a-R,a+R)$; continuity follows from uniform limits.
\item Since $T'(x)=1$, $f'(x)=g'(Tx)=\sum\limits_{n=1}^{\infty}n c_n(Tx)^{n-1}$, with local uniform convergence.
\item If $[y,z]\subset(a-R,a+R)$, choose $\rho<R$ with $[y,z]\subset[a-\rho,a+\rho]$.\\
Uniform convergence on $[a-\rho,a+\rho]$ gives $\int_y^z f=\lim\limits_N\int_y^z\sum\limits_{n=0}^{N}c_n(Tx)^n\,dx$, yielding the stated formula.
\end{enumerate}
\hlt{Structural Commentary}: Work on compact intervals strictly inside the radius.\\
\hlt{Key Insight}: Translation reduces centered series to ordinary power series.
\end{proof}
\end{corollary}

\begin{corollary} \label{cor:analytic-functions-smooth-and-unique-taylor} \hlt{Analytic Functions Are Their Taylor Series}\\
Suppose $f(x)=\sum\limits_{n=0}^{\infty}c_n(x-a)^n$ for $\abs{x-a}<r$, with $r>0$.
\begin{enumerate}[label=(\roman*)]
\setlength{\itemsep}{0pt}
\item $f\in C^\infty(a-r,a+r)$.
\item $c_k=f^{(k)}(a)/k!$ for every $k\geq0$.
\item The power-series expansion of $f$ about $a$ is unique.
\item $f(x)=\sum\limits_{n=0}^{\infty}\frac{f^{(n)}(a)}{n!}(x-a)^n$ for $\abs{x-a}<r$.
\end{enumerate}
\begin{proof}
\begin{enumerate}[label=(\roman*)]
\setlength{\itemsep}{0pt}
\item Termwise differentiation gives $f\in C^\infty(a-r,a+r)$.
\item For $k\geq0$, $f^{(k)}(x)=\sum\limits_{n\geq k}\frac{n!}{(n-k)!}c_n(x-a)^{n-k}$; at $x=a$, $f^{(k)}(a)=k!c_k$.
\item If $f=\sum b_n(x-a)^n$ also, then $b_k=f^{(k)}(a)/k!=c_k$ for all $k$.
\item Replace $c_n$ by $f^{(n)}(a)/n!$ in $f(x)=\sum c_n(x-a)^n$.
\end{enumerate}
\hlt{Structural Commentary}: Derivatives determine every coefficient.\\
\hlt{Key Insight}: Analytic representation forces its own Taylor series.
\end{proof}
\end{corollary}

\begin{theorem} \label{thm:power-series-reexpansion-and-identity} \hlt{Re-Expansion and Identity for Power Series}\\
Let $f(x)=\sum_{n\ge0}c_nx^n$ have radius $R>0$, and set $S=(-R,R)$.
\begin{enumerate}[label=(\roman*)]
\setlength{\itemsep}{0pt}
\item If $a\in S$, then $f(x)=\sum\limits_{m\geq0}\frac{f^{(m)}(a)}{m!}(x-a)^m$ for $\abs{x-a}<R-\abs a$.
\item Let $F(x)=\sum\limits_{n\geq0}a_nx^n$ and $G(x)=\sum\limits_{n\geq0}b_nx^n$ converge on $S$.\\
If $\mathcal{Z}_{F,G}\defeq\{x\in S:F(x)=G(x)\}$ has a limit point in $S$, then $a_n=b_n$ for all $n\ge0$, hence $F=G$ on $S$.
\end{enumerate}
\begin{proof}
\begin{enumerate}[label=(\roman*)]
\setlength{\itemsep}{0pt}
\item Put $\rho\defeq\abs{x-a}+\abs a<R$. Since $\sum\limits_{n\geq0}\abs{c_n}\rho^n<\infty$, the binomial double sum is absolutely rearrangeable:\\
$f(x)=\sum\limits_{m\geq0}A_m(x-a)^m$, where $A_m\defeq\sum\limits_{n\geq m}\binom{n}{m}c_na^{n-m}$. Also $f^{(m)}(a)=\sum\limits_{n\geq m}\frac{n!}{(n-m)!}c_na^{n-m}=m!A_m$.
\item Set $H\defeq F-G=\sum\limits_{n\geq0}d_nx^n$, $d_n\defeq a_n-b_n$, $Z\defeq\{x\in S:H(x)=0\}$, and $A\defeq Z'\cap S$.\\
For $x_0\in A$, (i) gives $H(x)=\sum e_n(x-x_0)^n$ near $x_0$.\\
If $k=\min\{n:e_n\neq0\}$ existed, then $H=(x-x_0)^kg$ with $g(x_0)=e_k\neq0$, so $x_0$ would be isolated in $Z$, contradiction. Thus all $e_n=0$, so $H=0$ near $x_0$ and $A$ is open. Since $S\setminus A$ is open, $A\neq\emptyset$, and $S$ is connected, $A=S$. As $Z$ is closed in $S$, $Z=S$; hence $d_n=H^{(n)}(0)/n!=0$.
\end{enumerate}
\hlt{Structural Commentary}: Re-expand locally; connectedness propagates zeros.\\
\hlt{Key Insight}: Accumulating zeros force every local coefficient to vanish.
\end{proof}
\end{theorem}

\begin{perspective} \label{per:regularity-hierarchy-one-variable} \hlt{Regularity Hierarchy}\\
The inclusions are $C^\omega\subseteq C^\infty\subseteq \cdots \subseteq C^2\subseteq C^1\subseteq C^0$.\\
The inclusions follow from Theorem \ref{thm:power-series-uniform-differentiation} and Definition \ref{def:smoothness-classes-one-variable}; they are generally proper.\\
\hlt{Key Insight}: Higher regularity is extra structure, not just a different notation for continuity.
\end{perspective}

\begin{example} \label{ex:smooth-flat-nonanalytic-function} \hlt{Smooth but Not Analytic}\\
Define $e(x)=
\begin{cases}
e^{-1/x},&x>0,\\
0,&x\leq0.
\end{cases}$.
By repeated L'Hospital, Theorem \ref{thm:lhospital-one-sided}, $e\in C^\infty$ and $e^{(r)}(0)=0$ for all $r$.\\
Its Taylor series at $0$ is zero, but $e(x)>0$ for $x>0$, so $e$ is not analytic at $0$ (Definition \ref{def:real-analytic-at-point}).\\
\hlt{Key Insight}: A function can be infinitely tangent to the $x$-axis without being locally equal to its Taylor series.
\end{example}

\begin{perspective} \label{per:taylor-coefficients-need-a-remainder} \hlt{Taylor Coefficients Are Only the First Half}\\
Power-series coefficients are forced by derivatives (Corollary \ref{cor:analytic-functions-smooth-and-unique-taylor}); for a function not already known to be analytic, proving the Taylor series represents the function needs a remainder estimate (Theorem \ref{thm:taylor-theorem-lagrange-remainder}).
\end{perspective}

\begin{theorem} \label{thm:exponential-logarithm-from-power-series} \hlt{Exponential and Logarithm from the Power Series}\\
Define $\mathsf{E}(z)\defeq\sum\limits_{n=0}^{\infty}\frac{z^n}{n!}$ for $z\in\C$.
\begin{enumerate}[label=(\roman*)]
\setlength{\itemsep}{0pt}
\item $\mathsf{E}(z+w)=\mathsf{E}(z)\mathsf{E}(w)$ for $z,w\in\C$.
\item On $\R$, $\mathsf{E}>0$, $\mathsf{E}'=\mathsf{E}$, $\mathsf{E}(x)=e^x$, and $\mathsf{E}$ is strictly increasing.
\item The inverse $\mathsf{L}\defeq\mathsf{E}^{-1}:(0,\infty)\to\R$ satisfies $\mathsf{L}'(y)=1/y$ and $\mathsf{L}(uv)=\mathsf{L}(u)+\mathsf{L}(v)$ for $u,v>0$.
\item For $x>0$ and $\alpha\in\R$, set $x^\alpha\defeq\mathsf{E}(\alpha\mathsf{L}(x))$; then $(x^\alpha)'=\alpha x^{\alpha-1}$.
\end{enumerate}
\begin{proof}
The ratio test gives convergence for all $z\in\C$.
\begin{enumerate}[label=(\roman*)]
\setlength{\itemsep}{0pt}
\item Absolute convergence gives $\mathsf{E}(z)\mathsf{E}(w)=\sum\limits_{n\geq0}\sum\limits_{j=0}^{n}\frac{z^jw^{n-j}}{j!(n-j)!}=\sum\limits_{n\geq0}\frac{(z+w)^n}{n!}=\mathsf{E}(z+w)$.
\item Termwise differentiation gives $\mathsf{E}'=\mathsf{E}$ on $\R$. By (i), $\mathsf{E}(x)\mathsf{E}(-x)=1$, so continuity and $\mathsf{E}(0)=1$ give $\mathsf{E}>0$; hence $\mathsf{E}'=\mathsf{E}>0$. Since $\mathsf{E}(1)=e$, (i) gives $\mathsf{E}(p)=e^p$ for $p\in\Q$, and continuity gives $\mathsf{E}(x)=e^x$.
\item For $y=\mathsf{E}(x)$, differentiating $\mathsf{L}(\mathsf{E}(x))=x$ gives $\mathsf{L}'(y)=1/\mathsf{E}'(x)=1/y$.\\
If $u=\mathsf{E}(s)$ and $v=\mathsf{E}(t)$, then $\mathsf{L}(uv)=\mathsf{L}(\mathsf{E}(s+t))=s+t=\mathsf{L}(u)+\mathsf{L}(v)$.
\item $\frac{d}{dx}\mathsf{E}(\alpha\mathsf{L}(x))=\mathsf{E}(\alpha\mathsf{L}(x))\alpha\mathsf{L}'(x)=x^\alpha\alpha/x=\alpha x^{\alpha-1}$.
\end{enumerate}
\hlt{Structural Commentary}: Addition, inversion, powers.\\
\hlt{Key Insight}: $\mathsf{E}(z+w)=\mathsf{E}(z)\mathsf{E}(w)$ creates the logarithm law.
\end{proof}
\end{theorem}

\begin{perspective} \label{per:real-powers-rational-approximation-route} \hlt{Real Powers by Rational Approximation}\\
One route defines $x^\alpha$ for $x>0$ by choosing $q_n\in\Q$ with $q_n\to\alpha$ and setting $x^\alpha=\lim_{n\to\infty}x^{q_n}$ after proving existence and independence of $(q_n)$. The power-series construction sets $x^\alpha=\mathsf{E}(\alpha\mathsf{L}(x))$; both routes must agree with rational powers and be choice-independent.
\end{perspective}

\begin{perspective} \label{per:ross-logarithmic-integral-route} \hlt{Logarithmic Integral Route}\\
With the Riemann integral, set $L(y)\defeq\int_1^y \frac1t\,dt$ for $y>0$. Theorem \ref{thm:integral-function-continuity-and-derivative} gives $L'(y)=1/y$; endpoint limits make $L$ invertible, and Theorem \ref{thm:chain-rule-one-variable} then recovers $E'=E$ and the product laws.
\end{perspective}

\begin{perspective} \label{per:exponential-factorial-and-improper-integrals} \hlt{From Exponentials to Factorials}\\
Theorems \ref{thm:exponential-logarithm-from-power-series}, \ref{thm:riemann-integration-by-parts} supply real-power laws and integration-by-parts pattern used for $\Gamma$ in Theorem \ref{thm:gamma-function-basic-theory}.
\end{perspective}

\begin{theorem} \label{thm:trigonometric-functions-from-exponential} \hlt{Trigonometric Functions from the Exponential}\\
For $x\in\R$, set $\mathsf{C}(x)\defeq\frac{\mathsf{E}(ix)+\mathsf{E}(-ix)}{2}$, $\mathsf{S}(x)\defeq\frac{\mathsf{E}(ix)-\mathsf{E}(-ix)}{2i}$, and $\gamma(t)\defeq\mathsf{E}(it)$.
\begin{enumerate}[label=(\roman*)]
\setlength{\itemsep}{0pt}
\item $\mathsf{E}(ix)=\mathsf{C}(x)+i\mathsf{S}(x)$, $\mathsf{C}^2+\mathsf{S}^2=1$, $\mathsf{C}'=-\mathsf{S}$, and $\mathsf{S}'=\mathsf{C}$.
\item $x_0\defeq\min\{x>0:\mathsf{C}(x)=0\}$ exists; with $\pi\defeq2x_0$, $\mathsf{E}(\pi i)=-1$, $\mathsf{E}(2\pi i)=1$, $\mathsf{E}(z+2\pi i)=\mathsf{E}(z)$ for $z\in\C$.
\item $\mathsf{C},\mathsf{S}$ are $2\pi$-periodic, $\gamma([0,2\pi])=\{\zeta\in\C:\abs{\zeta}=1\}$, and $\gamma$ is injective on $[0,2\pi)$.
\end{enumerate}
\begin{proof}
\begin{enumerate}[label=(\roman*)]
\setlength{\itemsep}{0pt}
\item Real coefficients give $\overline{\mathsf{E}(ix)}=\mathsf{E}(-ix)$, so $\mathsf{E}(ix)=\mathsf{C}+i\mathsf{S}$ and $\abs{\mathsf{E}(ix)}^2=\mathsf{E}(ix)\mathsf{E}(-ix)=1$.\\
Thus $\mathsf{C}^2+\mathsf{S}^2=1$; differentiating the definitions gives $\mathsf{C}'=-\mathsf{S}$ and $\mathsf{S}'=\mathsf{C}$.
\item If $\mathsf{C}$ has no positive zero, then $\mathsf{C}>0$ on $(0,\infty)$, hence $\mathsf{S}'>0$ and $\mathsf{S}(x)>0$ for $x>0$.\\
For $0<x<y$, $\mathsf{S}(x)(y-x)<\int_x^y\mathsf{S}(t)\,dt=\mathsf{C}(x)-\mathsf{C}(y)\leq2$, impossible as $y\to\infty$; continuity gives the first zero $x_0$. On $(0,x_0)$, $\mathsf{C}>0$, so $\mathsf{S}$ increases from $0$ to $\mathsf{S}(x_0)=1$.\\
Since $x_0=\pi/2$, $\mathsf{E}(\pi i/2)=i$, and $\mathsf{E}(z+w)=\mathsf{E}(z)\mathsf{E}(w)$ gives $\mathsf{E}(\pi i)=-1$, $\mathsf{E}(2\pi i)=1$, and $\mathsf{E}(z+2\pi i)=\mathsf{E}(z)$.
\item Periodicity of $\mathsf{C},\mathsf{S}$ follows from (ii). On $[0,\pi/2]$, $\mathsf{C}$ decreases from $1$ to $0$ and $\mathsf{S}$ increases from $0$ to $1$, so $\gamma$ traverses the first quadrant once. Since $\gamma(t+m\pi/2)=\gamma(t)i^m$ for $m=0,1,2,3$, the four quadrants give $\gamma([0,2\pi])=\{\zeta:\abs{\zeta}=1\}$ and injectivity on $[0,2\pi)$.
\end{enumerate}
\hlt{Structural Commentary}: $\mathsf{E}(it)$ supplies trigonometry.\\
\hlt{Key Insight}: $\mathsf{E}(it)\mathsf{E}(-it)=1$ puts the path on the unit circle.
\end{proof}
\end{theorem}

\begin{figure}[H]
  \centering
  \begin{tikzpicture}[
  scale=1.25,
  >=Stealth,
  every node/.style={font=\scriptsize},
  axis/.style={->, gray!70, line width=0.45pt},
  circlecurve/.style={blue!70!black, line width=0.9pt},
  arcpath/.style={orange!80!black, line width=0.9pt},
  guide/.style={gray!45, dashed, line width=0.4pt},
  point/.style={circle, fill=blue!70!black, inner sep=1.4pt},
  labelbox/.style={fill=white, fill opacity=0.9, text opacity=1, inner sep=1.2pt}
]
  \draw[axis] (-1.35,0) -- (1.5,0) node[right] {$\operatorname{Re}$};
  \draw[axis] (0,-1.35) -- (0,1.5) node[above] {$\operatorname{Im}$};

  \draw[circlecurve] (0,0) circle (1);
  \draw[arcpath, ->] (1,0) arc[start angle=0,end angle=55,radius=1];

  \coordinate (P) at (55:1);
  \draw[guide] (P) -- ({cos(55)},0) node[below] {$\mathsf{C}(t)$};
  \draw[guide] (P) -- (0,{sin(55)}) node[left] {$\mathsf{S}(t)$};
  \draw[guide] (0,0) -- (P);

  \node[point] at (P) {};
  \node[labelbox, anchor=south west] at (P) {$\mathsf{E}(it)=\mathsf{C}(t)+i\mathsf{S}(t)$};
  \node[labelbox] at (0.5,0.18) {$t$};
  \node[below right] at (1,0) {$1$};
  \node[above left] at (0,1) {$i$};
\end{tikzpicture}

  \caption{$\mathsf{E}(it)=\mathsf{C}(t)+i\mathsf{S}(t)$ traces the unit circle, defining sine and cosine analytically.}
  \label{fig:exponential-unit-circle}
\end{figure}

\begin{theorem} \label{thm:fundamental-theorem-algebra-exponential-proof} \hlt{Algebraic Completeness of $\C$}\\
If $P\in\C[z]$ and $\deg P\geq1$, then $\exists z_0\in\C$ such that $P(z_0)=0$.
\begin{proof} Write $P(z)=\sum_{j=0}^{n}a_jz^j$, $n\geq1$, $a_n=1$.\\
Since $\abs{P(z)}\to\infty$ as $\abs z\to\infty$, choose $z_0$ with $\mu\defeq\abs{P(z_0)}=\min_{\C}\abs P$. If $\mu=0$, done.\\
Assume $\mu>0$ and set $Q(w)\defeq P(z_0+w)/P(z_0)$. Then $Q(0)=1$ and $\abs{Q(w)}\geq1$ for all $w$.\\
Write $Q(w)=1+b_kw^k+w^{k+1}H(w)$ with $k\geq1$, $b_k\neq0$, and $H$ polynomial.\\
Choose $\theta\in\R$ with $b_ke^{ik\theta}=-\abs{b_k}$. For $r>0$, $Q(re^{i\theta})=1-r^kA(r)$, where $A(r)\defeq\abs{b_k}-rH(re^{i\theta})$.\\
For small $r$, $\Re A(r)>\abs{b_k}/2$, hence $\abs{1-r^kA(r)}^2=1-2r^k\Re A(r)+r^{2k}\abs{A(r)}^2<1$, contradicting $\abs{Q}\geq1$.\\
\hlt{Structural Commentary}: Minimize $\abs P$, translate, perturb.\\
\hlt{Key Insight}: A phase choice makes the first nonzero term decrease $\abs Q$.
\end{proof}
\end{theorem}

\begin{theorem} \label{thm:gamma-function-basic-theory} \hlt{Gamma Function, Beta Function, and Stirling Skeleton}\\
For $x>0$, set $\Gamma(x)\defeq\int_0^\infty t^{x-1}e^{-t}\,dt$; for $x,y>0$, set $\mathrm{B}(x,y)\defeq\int_0^1t^{x-1}(1-t)^{y-1}\,dt$.
\begin{enumerate}[label=(\roman*)]
\setlength{\itemsep}{0pt}
\item $\Gamma(x+1)=x\Gamma(x)$ for $x>0$, $\Gamma(n+1)=n!$ for $n\in\N$, and $\log\Gamma$ is convex on $(0,\infty)$.
\item $\Gamma$ is the unique $f:(0,\infty)\to(0,\infty)$ with $f(1)=1$, $f(x+1)=xf(x)$, and $\log f$ convex.
\item $\mathrm{B}(x,y)=\frac{\Gamma(x)\Gamma(y)}{\Gamma(x+y)}$ for $x,y>0$.
\item $\Gamma\!\left(\frac12\right)=\sqrt{\pi}$, $\int_{-\infty}^{\infty}e^{-s^2}\,ds=\sqrt{\pi}$, $\Gamma(x)=\frac{2^{x-1}}{\sqrt{\pi}}\Gamma\!\left(\frac{x}{2}\right)\Gamma\!\left(\frac{x+1}{2}\right)$, and $\lim\limits_{x\to\infty}\frac{\Gamma(x+1)}{(x/e)^x\sqrt{2\pi x}}=1$.
\end{enumerate}
\begin{proof}
\begin{enumerate}[label=(\roman*)]
\setlength{\itemsep}{0pt}
\item Integration by parts gives $\Gamma(x+1)=\int_0^\infty t^xe^{-t}\,dt=x\Gamma(x)$; $\Gamma(1)=1$ gives $\Gamma(n+1)=n!$.\\
For $0<\lambda<1$, Holder gives $\Gamma(\lambda x+(1-\lambda)y)\leq\Gamma(x)^\lambda\Gamma(y)^{1-\lambda}$, hence $\log\Gamma$ is convex.
\item Let $\varphi=\log f$. For $0<x<1$ and $n\geq1$, convexity with $\varphi(t+1)=\varphi(t)+\log t$ gives\\
$0\leq\varphi(x)-\log\!\frac{n!\,n^x}{x(x+1)\cdots(x+n)}\leq x\log(1+\frac1n)$. Letting $n\to\infty$ fixes $\varphi$ on $(0,1)$; recurrence fixes $\varphi$ on $(0,\infty)$.
\item Fix $y>0$ and set $F_y(x)\defeq\frac{\Gamma(x+y)}{\Gamma(y)}\mathrm{B}(x,y)$. Since $\mathrm{B}(1,y)=1/y$ and $\mathrm{B}(x+1,y)=\frac{x}{x+y}\mathrm{B}(x,y)$, $F_y(1)=1$ and $F_y(x+1)=xF_y(x)$; Holder gives log-convexity. Thus (ii) gives $F_y=\Gamma$.
\item $t=\sin^2\theta$ gives $\mathrm{B}(\frac12,\frac12)=\pi$, hence $\Gamma(\frac12)=\sqrt{\pi}$ by (iii). Then $t=s^2$ gives $\Gamma(\frac12)=\int_{-\infty}^{\infty}e^{-s^2}\,ds$.\\
Beta duplication gives the duplication formula, and $t=x(1+u)$ with Laplace asymptotics gives Stirling.
\end{enumerate}
\hlt{Structural Commentary}: Recurrence plus log-convexity gives uniqueness.\\
\hlt{Key Insight}: Bohr--Mollerup turns beta identities into gamma identities.
\end{proof}
\end{theorem}

\begin{method} \label{meth:gamma-function-proof-pattern} \hlt{Gamma Function Proof Pattern}
\begin{enumerate}[label=(\roman*)]
\setlength{\itemsep}{0pt}
\item \hlt{Gamma function}: use $\Gamma(x+1)=x\Gamma(x)$, $\Gamma(1)=1$, and log-convexity for uniqueness.
\end{enumerate}
\end{method}
